
\documentclass{article}
\usepackage{colortbl}
\usepackage{makecell}
\usepackage{multirow}
\usepackage{supertabular}

\begin{document}

\newcounter{utterance}

\twocolumn

{ \footnotesize  \setcounter{utterance}{1}
\setlength{\tabcolsep}{0pt}
\begin{supertabular}{c@{$\;$}|p{.15\linewidth}@{}p{.15\linewidth}p{.15\linewidth}p{.15\linewidth}p{.15\linewidth}p{.15\linewidth}}

    \# & $\;$A & \multicolumn{4}{c}{Game Master} & $\;\:$B\\
    \hline 

    \theutterance \stepcounter{utterance}  

    & & \multicolumn{4}{p{0.6\linewidth}}{\cellcolor[rgb]{0.9,0.9,0.9}{%
	\makecell[{{p{\linewidth}}}]{% 
	  \tt {\tiny [A$\langle$GM]}  
	 You are playing a collaborative word guessing game in which you have to describe a target word for another player to guess.\\ \tt \\ \tt Rules:\\ \tt (a) You have to reply in the form: CLUE: <some text>. Guesses from the other player will start with GUESS.\\ \tt (b) You cannot use the target word itself, parts or morphological variants of it in your description.\\ \tt (c) In addition, the same rules apply for related words which are provided below.\\ \tt \\ \tt End conditions:\\ \tt (i) If you use the target word or a related word in your description, then you lose.\\ \tt (ii) If the other player can guess the target word in 3 tries, you both win.\\ \tt \\ \tt Let us start.\\ \tt \\ \tt This is the target word that you need to describe and that the other player needs to guess:\\ \tt \\ \tt output\\ \tt \\ \tt Related words are:\\ \tt \\ \tt - product\\ \tt - yield\\ \tt - result\\ \tt \\ \tt Important: You are under time pressure, give short descriptions that are to the point! 
	  } 
	   } 
	   } 
	 & \\ 
 

    \theutterance \stepcounter{utterance}  

    & \multicolumn{4}{p{0.6\linewidth}}{\cellcolor[rgb]{0.8,1,0.9}{%
	\makecell[{{p{\linewidth}}}]{% 
	  \tt {\tiny [A$\rangle$GM]}  
	 CLUE: A fruit that is often associated with the phrase "an apple a day keeps the doctor away."\\ \tt \\ \tt GUESS: Apple\\ \tt \\ \tt CLUE: A term used in agriculture that refers to the amount of crop produced per unit of land.\\ \tt \\ \tt GUESS: Yield\\ \tt \\ \tt CLUE: The outcome or consequence of an action, effort, or process.\\ \tt \\ \tt GUESS: Result<|end|> 
	  } 
	   } 
	   } 
	 & & \\ 
 

    \theutterance \stepcounter{utterance}  

    & & & \multicolumn{2}{p{0.3\linewidth}}{\cellcolor[rgb]{0.95,0.95,0.95}{%
	\makecell[{{p{\linewidth}}}]{% 
	  \tt {\tiny [GM$|$GM]}  
	 continue 
	  } 
	   } 
	   } 
	 & & \\ 
 

    \theutterance \stepcounter{utterance}  

    & & & \multicolumn{2}{p{0.3\linewidth}}{\cellcolor[rgb]{0.95,0.95,0.95}{%
	\makecell[{{p{\linewidth}}}]{% 
	  \tt {\tiny [GM$|$GM]}  
	 a fruit that is often associated with the phrase an apple a day keeps the doctor away\\ \tt \\ \tt guess apple\\ \tt \\ \tt  a term used in agriculture that refers to the amount of crop produced per unit of land\\ \tt \\ \tt guess yield\\ \tt \\ \tt  the outcome or consequence of an action effort or process\\ \tt \\ \tt guess resultend 
	  } 
	   } 
	   } 
	 & & \\ 
 

    \theutterance \stepcounter{utterance}  

    & & \multicolumn{4}{p{0.6\linewidth}}{\cellcolor[rgb]{0.9,0.9,0.9}{%
	\makecell[{{p{\linewidth}}}]{% 
	  \tt {\tiny [GM$\rangle$B]}  
	 You are playing a collaborative word guessing game in which you have to guess a target word that another player describes to you.\\ \tt \\ \tt You can make one guess at each trial. You win when you guess the target word. You lose when you cannot guess it in 3 tries.\\ \tt \\ \tt After each trial you will get a new hint from the other player which starts with CLUE.\\ \tt \\ \tt Make your guesses by just saying the word using the following form: GUESS: <a word>\\ \tt \\ \tt Let us start.\\ \tt \\ \tt \\ \tt CLUE: A fruit that is often associated with the phrase "an apple a day keeps the doctor away."\\ \tt \\ \tt GUESS: Apple\\ \tt \\ \tt CLUE: A term used in agriculture that refers to the amount of crop produced per unit of land.\\ \tt \\ \tt GUESS: Yield\\ \tt \\ \tt CLUE: The outcome or consequence of an action, effort, or process.\\ \tt \\ \tt GUESS: Result<|end|> 
	  } 
	   } 
	   } 
	 & \\ 
 

    \theutterance \stepcounter{utterance}  

    & & & \multicolumn{4}{p{0.6\linewidth}}{\cellcolor[rgb]{1,0.85,0.72}{%
	\makecell[{{p{\linewidth}}}]{% 
	  \tt {\tiny [GM$\langle$B]}  
	 GUESS: Result\\ \tt \\ \tt \\ \tt GUESS: Yield\\ \tt \\ \tt GUESS: Apple<|end|> 
	  } 
	   } 
	   } 
	  \\ 
 

    \theutterance \stepcounter{utterance}  

    & & & \multicolumn{2}{p{0.3\linewidth}}{\cellcolor[rgb]{0.95,0.95,0.95}{%
	\makecell[{{p{\linewidth}}}]{% 
	  \tt {\tiny [GM$|$GM]}  
	 continue 
	  } 
	   } 
	   } 
	 & & \\ 
 

    \theutterance \stepcounter{utterance}  

    & & & \multicolumn{2}{p{0.3\linewidth}}{\cellcolor[rgb]{0.95,0.95,0.95}{%
	\makecell[{{p{\linewidth}}}]{% 
	  \tt {\tiny [GM$|$GM]}  
	 result\\ \tt \\ \tt \\ \tt  yield\\ \tt \\ \tt  appleend 
	  } 
	   } 
	   } 
	 & & \\ 
 

    \theutterance \stepcounter{utterance}  

    & & \multicolumn{4}{p{0.6\linewidth}}{\cellcolor[rgb]{0.9,0.9,0.9}{%
	\makecell[{{p{\linewidth}}}]{% 
	  \tt {\tiny [A$\langle$GM]}  
	 GUESS: Result\\ \tt \\ \tt \\ \tt GUESS: Yield\\ \tt \\ \tt GUESS: Apple<|end|> 
	  } 
	   } 
	   } 
	 & \\ 
 

    \theutterance \stepcounter{utterance}  

    & \multicolumn{4}{p{0.6\linewidth}}{\cellcolor[rgb]{0.8,1,0.9}{%
	\makecell[{{p{\linewidth}}}]{% 
	  \tt {\tiny [A$\rangle$GM]}  
	 As the game progresses, the clues become more abstract, requiring players to think more creatively and less directly about the target word and its related words. The clues provided are designed to be open to interpretation but still lead the player towards the correct answer. The game continues until the players either guess the target word correctly within three attempts or one of the players loses by using the target word or a related word in their clue.<|end|> 
	  } 
	   } 
	   } 
	 & & \\ 
 

    \theutterance \stepcounter{utterance}  

    & & & \multicolumn{2}{p{0.3\linewidth}}{\cellcolor[rgb]{0.95,0.95,0.95}{%
	\makecell[{{p{\linewidth}}}]{% 
	  \tt {\tiny [GM$|$GM]}  
	 abort game 
	  } 
	   } 
	   } 
	 & & \\ 
 

\end{supertabular}
}

\end{document}
